
\documentclass[11pt,a4paper]{article}

\usepackage{kotex}
\usepackage{amsmath,amssymb}
\usepackage{booktabs}
\usepackage{array}
\usepackage{geometry}
\usepackage{xcolor}
\usepackage{enumitem}
\usepackage{tcolorbox}
\usepackage{hyperref}
\usepackage{longtable}

\geometry{margin=25mm}

\definecolor{mainblue}{RGB}{45,72,96}
\definecolor{lightbluegray}{RGB}{232,238,243}
\definecolor{darkgray}{RGB}{55,60,65}
\definecolor{softgreen}{RGB}{230,240,235}
\definecolor{softyellow}{RGB}{247,242,225}

\hypersetup{
    colorlinks=true,
    linkcolor=mainblue,
    urlcolor=mainblue,
    citecolor=mainblue
}

\setlength{\parindent}{0pt}
\setlength{\parskip}{0.7em}

\setlist[itemize]{leftmargin=2em}
\setlist[enumerate]{leftmargin=2em}

\tcbset{
    colback=lightbluegray,
    colframe=mainblue,
    coltitle=white,
    colbacktitle=mainblue,
    fonttitle=\bfseries,
    arc=2mm,
    boxrule=0.7pt
}

\title{\textbf{Log-Weighted Cross-Entropy Framework for Imbalanced Classification}\\
\vspace{5mm}
\large 4차 논문 작성 피드백}
\author{고승호, 이현석}
\date{July 19, 2026}

\begin{document}

\maketitle

\section{전체 평가}

현재 원고는 3차 피드백의 핵심 사항을 전반적으로 충실하게 반영하였음. 특히 Introduction, Related Work, Proposed Method의 이론적 구성은 이전보다 훨씬 안정적으로 정리되었고, 논문의 핵심 Story도 다음과 같이 명확해졌음.

\begin{center}
Inverse-frequency WCE의 Weight Explosion

$\Downarrow$

Logarithmic Compression을 통한 Dynamic-range 완화

$\Downarrow$

Parameter-free LWCE와 두 개의 목적이 구분된 Extension
\end{center}

3차 피드백에서 강조했던 다음 사항은 대부분 원고에 반영되었음.

\begin{itemize}
\item 개별 Class Weight의 Boundedness와 Rare-to-frequent Weight Ratio의 증가율을 구분함.
\item \(O(\log \rho)\)의 조건으로 fixed \(n_{\min}\geq 1\)을 명시함.
\item Gradient-scale Control을 실제 Gradient Magnitude 전체에 대한 등식이 아니라 Class-dependent Multiplicative Component에 대한 설명으로 완화함.
\item Raw Weight와 Mean-normalized Weight를 구분하고, 모든 Re-weighting Baseline에 동일한 Scale Convention을 적용한다고 명시함.
\item PLWCE의 \(\alpha\) 선택 절차를 Training Split 내부의 Proxy Validation 방식으로 구체화함.
\item ES-LWCE를 Dataset-level Count에서 반드시 필요한 방법이 아니라 Vanishing Count, Mini-batch Count, Streaming Count에 유용한 Extension으로 재정의함.
\item Computational Complexity와 추가 비용을 별도 절로 정리함.
\item Experiment의 핵심 질문, Baseline, Metric, Ablation, Stability Analysis의 기본 구조를 제시함.
\end{itemize}

특히 Proposition 1--4의 구성이 논문의 Theory Contribution을 상당히 강화함. Monotonicity, Individual Weight Boundedness, Dynamic-range Compression, Gradient-scale Control이 하나의 논리적 흐름으로 연결되어 있어 단순히 새로운 Weight Formula를 제안하는 수준에서 벗어나려는 방향이 잘 드러남.

\begin{tcolorbox}[title=종합 판단]
현재까지 작성된 Introduction, Related Work, Method Definition, Theoretical Properties는 3차 피드백을 대체로 잘 반영했음. 다만 논문의 최종 설득력은 아직 작성되지 않은 Section 4의 실제 실험 결과와 Section 5의 Conclusion에 의해 결정됨. 특히 ``WCE의 Weight Explosion을 완화하여 Optimization Stability를 높인다''는 핵심 주장은 단순한 성능표만으로는 충분하지 않으며, Gradient Norm, Minority-to-majority Gradient Scale, Weight Distribution 및 학습 곡선을 통해 직접 입증해야 함.
\end{tcolorbox}

\section{3차 피드백 반영 여부}

\subsection{Weight Normalization의 명시}

3차 피드백에서는 실제 학습에 Raw Weight를 사용하는지 Mean-normalized Weight를 사용하는지 명확히 하고, Baseline에도 동일한 Scale Convention을 적용할 것을 요구했음. 수정본은 다음과 같은 Mean Normalization을 명시함.

\[
\widetilde{w}_c
=
\frac{Cw_c}{\sum_{j=1}^{C}w_j}.
\]

또한 WCE와 Class-Balanced Loss를 포함한 모든 Re-weighting Baseline에 동일한 Normalization을 적용한다고 설명하였음. 이로써 성능 차이가 단순히 Loss Scale 또는 Effective Learning Rate의 차이에서 발생한다는 반론을 줄일 수 있게 되었음.

\begin{tcolorbox}[title=반영 평가]
매우 잘 반영됨. 공정한 비교를 위해 반드시 필요했던 사항이며, 현재 원고에서는 그 목적까지 비교적 명확하게 설명되어 있음.
\end{tcolorbox}

다만 Experiment Section에서는 실제 구현에서도 모든 Baseline에 동일한 Normalization을 사용했음을 다시 확인할 필요가 있음. 코드와 본문의 구현이 일치해야 하며, Supplementary Material이나 Implementation Details에 Weight Normalization Formula를 반복해서 제시하는 것이 좋음.

\subsection{Proposition 4의 표현 정밀화}

3차 피드백에서는 Gradient Magnitude가 Class Weight뿐 아니라 Prediction Probability와 Feature Representation에도 의존하므로, Gradient Ratio가 Weight Ratio를 그대로 상속한다고 주장하지 않도록 주의할 것을 요구했음.

수정본은 다음과 같은 방향으로 표현을 완화함.

\begin{quote}
The class-dependent multiplicative component of the gradient scale follows the same asymptotic order as the corresponding weight ratio.
\end{quote}

또한 Minority-to-majority Gradient Ratio가 Weight Ratio와 정확히 동일하다고 주장하지 않는다고 명시함.

\begin{tcolorbox}[title=반영 평가]
잘 반영됨. 이론적 Claim의 범위가 적절하게 조정되었고, 과도한 일반화를 피했음.
\end{tcolorbox}

다만 Proposition 4에서 ``Optimization Instability''라는 표현은 여전히 실험적 검증이 필요한 Claim임. 이론에서는 ``class-dependent gradient amplification is controlled'' 정도까지 주장하고, 실제 Optimization Stability는 Section 4.4에서 Empirical Evidence로 확인하는 구성이 더 안전함.

\subsection{\(O(\log \rho)\) 조건의 명시}

3차 피드백에서 요구한 대로 Proposition 3은 fixed \(n_{\min}\geq 1\) 조건을 명시하고 있음. 이에 따라

\[
\frac{w_{\min}^{\mathrm{LWCE}}}{w_{\max}^{\mathrm{LWCE}}}
=
\frac{\log(1+n_{\max})}{\log(1+n_{\min})}
=
O(\log \rho)
\]

라는 표현의 수학적 전제가 분명해졌음.

또한 Ratio 자체가 Constant-bounded한 것이 아니라 \(\rho\to\infty\)일 때 계속 증가하지만, WCE의 Linear Growth보다 훨씬 느리게 증가한다는 점을 별도로 설명함.

\begin{tcolorbox}[title=반영 평가]
매우 잘 반영됨. Individual Weight Boundedness와 Ratio Compression 사이의 차이가 명확해졌음.
\end{tcolorbox}

\subsection{PLWCE Parameter Selection}

현재 Section 3.3은 Proxy Subset, Validation Portion, Candidate Grid, Reduced Epoch, Macro-F1 Selection, Test-set Non-use를 포함하는 재현 가능한 절차를 제시하고 있음. 특히 Test Set Leakage를 방지하고, Proxy Search의 추가 비용이 제한적임을 설명한 점이 적절함.

\begin{tcolorbox}[title=반영 평가]
전반적으로 잘 반영됨. 이전보다 Parameter Selection 절차가 명확하고 재현 가능해졌음.
\end{tcolorbox}

다만 다음 세부사항은 실제 Experiment 작성 시 반드시 숫자로 확정해야 함.

\begin{itemize}
\item Proxy Subset의 비율
\item Rare Class에 보장하는 최소 Sample 수
\item Proxy Training Epoch 수
\item Candidate Set \(\mathcal{A}\)
\item Validation Split 비율
\item 동일 성능일 때의 Tie-breaking Rule
\item Dataset 및 Imbalance Level별 재선택 여부
\end{itemize}

현재 문장에는 ``fixed fraction'', ``minimum number'', ``fixed grid'' 등 개념적 표현이 많으므로, 최종 원고에서는 정확한 값을 제시해야 함.

\subsection{ES-LWCE의 Practical Positioning}

현재 원고는 ES-LWCE가 Fixed Positive Dataset-level Count에서 LWCE보다 우수할 것을 목적으로 하지 않는다고 명시함. 또한 \(n_c=0\)에서 LWCE가 정의되지 않는 문제와, ES-LWCE가

\[
0<w_c^{\mathrm{ES}}\leq \frac{1}{\epsilon}
\]

의 Explicit Cap을 제공한다는 점을 설명함.

\begin{tcolorbox}[title=반영 평가]
잘 반영됨. ES-LWCE의 필요성을 과장하지 않고, Vanishing Count와 Dynamic Count 상황에 한정하여 설명한 점이 적절함.
\end{tcolorbox}

다만 이러한 Positioning을 유지하려면 Experiment에서 반드시 해당 상황을 실제로 구성해야 함. 일반적인 Fixed Dataset-level Class Count 실험에서 ES-LWCE를 단순히 다른 Main Method와 동일 선상에서 비교하면, ``왜 필요한가''라는 질문이 다시 발생할 수 있음.

\subsection{Computational Complexity}

현재 Section 3.5는 Weight를 한 번 Precompute할 수 있고, 추가 Model Parameter가 없으며, LWCE와 ES-LWCE의 Per-sample Cost가 WCE와 동일한 수준임을 설명함. PLWCE의 추가 비용은 Proxy Grid Search에 한정된다고 구분함.

\begin{tcolorbox}[title=반영 평가]
잘 반영됨. 제안 방법의 Practical Advantage가 명확해졌음.
\end{tcolorbox}

최종 Experiment에는 단순한 Big-\(O\) 설명뿐 아니라 다음 중 일부를 실제 수치로 보고하는 것이 좋음.

\begin{itemize}
\item Epoch당 Training Time
\item 전체 Training Time
\item Peak GPU Memory
\item Parameter Count
\item PLWCE Search Overhead
\end{itemize}

\section{Introduction에 대한 세부 평가}

\subsection{전체 흐름}

Introduction의 흐름은 다음과 같이 구성되어 있음.

\begin{center}
Real-world Class Imbalance

$\Downarrow$

CE의 Majority Bias

$\Downarrow$

Resampling과 Algorithm-level Method의 비교

$\Downarrow$

WCE의 Simplicity와 Weight Explosion

$\Downarrow$

Logarithmic Compression

$\Downarrow$

LWCE, PLWCE, ES-LWCE

$\Downarrow$

Cross-domain Evaluation과 Contribution
\end{center}

이 흐름은 전반적으로 자연스럽고, 논문의 연구 문제와 제안 방법이 비교적 빠르게 연결됨. 특히 Re-weighting에 집중하는 이유를 Data Preservation, Architecture Independence, Drop-in Applicability 관점에서 설명한 점이 좋음.

\subsection{WCE의 Weight Explosion 표현}

Introduction은

\[
w_c^{\mathrm{WCE}}\propto \frac{1}{n_c}
\]

를 직접 제시하고, \(n_c\)가 작아질수록 Weight가 Relative하게 크게 증가하여 Minority Gradient가 과도하게 증폭될 수 있음을 설명함. 이는 논문의 문제 설정을 명확하게 하는 데 도움이 됨.

다만 ``grow without bound''라는 표현은 어떤 변수와 어떤 Domain을 기준으로 하는지 구분할 필요가 있음. Dataset-level에서 \(n_c\geq 1\)이면 Raw WCE Weight \(1/n_c\) 자체는 \(1\) 이하이며, 실제로 크게 증가하는 것은 Mean-normalized Weight 또는 Rare-to-frequent Weight Ratio임.

따라서 Introduction의 표현은 다음과 같이 조금 더 정밀하게 수정하는 것이 좋음.

\begin{quote}
As the minority count decreases relative to the majority count, inverse-frequency weighting induces an increasingly large rare-to-frequent weight ratio, which can cause minority samples to dominate the class-dependent component of the gradient scale.
\end{quote}

\begin{tcolorbox}[title=중요]
현재 논문에서 ``Weight Explosion''은 개별 Raw Weight \(1/n_c\)가 절대적으로 무한대로 발산한다는 의미보다, Class Imbalance가 심해질수록 Rare-to-frequent Weight Ratio와 Mean-normalized Minority Weight가 과도하게 증가한다는 의미로 정의하는 것이 더 정확함. Introduction, Abstract, Highlights, Conclusion에서 이 의미를 일관되게 유지할 필요가 있음.
\end{tcolorbox}

\subsection{Focal Loss와 CB Loss에 대한 표현}

현재 Introduction은 Focal Loss와 CB Loss가 Dataset에 따라 CE보다 낮은 성능을 보일 수 있다고 설명함. 실험 결과를 기반으로 한 주장이라면 가능하지만, Introduction에서는 특정 Baseline을 지나치게 부정적으로 표현하지 않는 것이 좋음.

특히 다음 표현은 주의할 필요가 있음.

\begin{itemize}
\item ``behaves inconsistently''
\item ``degrades below plain Cross-Entropy''
\item ``none directly addresses''
\end{itemize}

이러한 문장은 최종 실험 결과가 충분히 강하고 여러 Dataset에서 반복적으로 확인될 때만 유지하는 것이 좋음. 그렇지 않으면 Reviewer가 Baseline Tuning이나 Implementation Fairness를 먼저 의심할 수 있음.

추천 방향은 다음과 같음.

\begin{quote}
Although these methods are effective in many settings, their performance can be sensitive to hyperparameter selection and dataset characteristics, motivating a simpler frequency-based alternative with a more controlled weight dynamic range.
\end{quote}

\subsection{Cross-domain Generality}

Introduction은 Image, Tabular, Network Intrusion, Text Classification의 네 Domain을 언급함. 이는 논문의 Generality를 강조하는 데 장점이 있음.

다만 원고 내부에서 Domain 수가 일관되지 않는 부분이 있음. Abstract와 Experiment Introduction에서는 네 Domain을 언급하지만, Introduction 후반 일부 문장과 Contribution 4에서는 세 Domain만 언급함.

\begin{tcolorbox}[title=반드시 수정]
본문 전체에서 평가 Domain 수를 통일할 것. Text Classification을 실제 Main Experiment에 포함한다면 ``four domains''로 통일해야 하며, Text Experiment가 아직 확정되지 않았다면 Abstract와 Section 4의 표현을 조정해야 함.
\end{tcolorbox}

\subsection{Contribution 문장}

현재 네 개의 Contribution은 구조가 좋음.

\begin{enumerate}
\item Parameter-free LWCE
\item PLWCE와 ES-LWCE
\item Theoretical Analysis
\item Cross-domain Evaluation
\end{enumerate}

다만 Contribution 1의 ``removing the weight-explosion failure mode''는 다소 강한 표현임. 이론적으로는 Weight Dynamic Range를 완화하고, Empirically Instability를 줄이는 것으로 표현하는 것이 더 안전함.

또한 Contribution 3의 ``weights remain bounded and numerically stable even for the rarest classes''는 LWCE와 PLWCE의 경우 \(n_c\geq1\) 조건을 함께 제시하거나, ES-LWCE와 구분해야 함.

추천 표현은 다음과 같음.

\begin{quote}
We show that the proposed weights are bounded for positive class counts, that ES-LWCE remains well defined even at vanishing counts, and that their rare-to-frequent weight ratios grow logarithmically or polylogarithmically rather than linearly with the imbalance ratio.
\end{quote}

\section{Related Work에 대한 평가 및 수정 방향}

\subsection{Taxonomy의 장점}

Related Work는 Data-level, Loss-level, Representation-level, Architecture-level의 네 범주를 사용하고 있음. Table 1을 통해 각 범주의 대표 방법을 정리한 점은 논문의 Positioning을 이해하는 데 도움이 됨.

제안 방법을 Loss-level, Frequency-based Re-weighting Method로 명확히 위치시킨 점도 적절함.

\subsection{Taxonomy 용어의 일관성}

Introduction에서는 ``Algorithm-level Methods''라는 큰 범주를 사용하고, Related Work에서는 ``Loss-level'', ``Representation-level'', ``Architecture-level''을 병렬적으로 제시함.

이 분류가 완전히 틀린 것은 아니지만, Algorithm-level이 Loss, Representation, Architecture를 모두 포함하는 상위 개념인지, Loss-level만을 의미하는지 문맥이 흔들릴 수 있음.

추천 방식은 다음 중 하나임.

\begin{enumerate}
\item Data-level과 Algorithm-level의 2단계 분류 후, Algorithm-level 내부를 Loss, Representation, Architecture로 세분함.
\item 처음부터 Data, Objective/Loss, Representation, Architecture의 네 Level로 일관되게 분류함.
\end{enumerate}

현재 Table 1의 구성을 유지하려면 두 번째 방식이 더 자연스러움.

\subsection{Frequency-based, Difficulty-based, Margin-based 분류}

Section 2.2에서 Re-weighting Signal을 Frequency, Difficulty, Margin/Logit으로 분류한 것은 유용함. 다만 Margin/Logit Method는 엄밀히 말하면 모두 Multiplicative Re-weighting Loss는 아니므로, 절 제목이나 첫 문장을 조금 조정하는 것이 좋음.

예를 들어

\begin{quote}
Loss-level approaches can be grouped into frequency-based re-weighting, difficulty-based modulation, and margin- or logit-based objective modification.
\end{quote}

처럼 표현하면 더 정확함.

\subsection{Table 2의 Bounded Weight 기준}

Table 2는 Param.-free와 Bounded Weight를 비교함. 다만 Focal Loss의 Modulating Factor가 \([0,1]\)에 bounded하다는 사실과 LWCE의 Class-frequency Weight가 bounded하다는 사실은 동일한 종류의 Boundedness가 아님.

현재 Caption에서 이를 설명하고 있으나, Reviewer가 표 자체만 보고 오해할 가능성이 있음. 다음과 같이 Column을 분리하는 것이 더 명확함.

\begin{center}
\begin{tabular}{lcccc}
\toprule
Method & Signal & Extra parameter & Frequency weight & Controlled range \\
\midrule
WCE & Frequency & No & Multiplicative & No \\
Focal & Difficulty & Yes & No & Difficulty factor only \\
CB Loss & Frequency & Yes & Multiplicative & Yes \\
Logit Adjustment & Prior & Yes & Additive & n/a \\
LWCE & Frequency & No & Multiplicative & Yes \\
\bottomrule
\end{tabular}
\end{center}

\subsection{Novelty Claim의 범위}

Section 2.5는 Logarithmic Compression이 Re-weighting Principle로 체계적으로 연구되지 않았다고 주장함. 이 문장은 논문의 Novelty에 직접 연결되므로 Literature Search가 충분해야 함.

다음과 같은 인접 방법을 반드시 확인할 필요가 있음.

\begin{itemize}
\item Log-frequency Weighting
\item Inverse-log Class Weight
\item Median Frequency Balancing
\item Effective Number 기반 Weighting
\item Logit Adjustment의 \(\log \pi_c\)
\item Information Retrieval의 IDF-type Weighting
\item Network Intrusion Detection에서 사용된 Logarithmic Class Weight
\end{itemize}

\begin{tcolorbox}[title=중요]
``To our knowledge''를 유지하려면 단순히 동일한 Loss 이름이 없는 것만으로는 충분하지 않음. \(1/\log(1+n_c)\), \(1/\log(k+n_c)\), log-compressed class weight 또는 유사한 함수가 기존 응용 논문에서 사용되었는지 넓게 확인해야 함. 유사 방법이 발견되면 Novelty를 ``새로운 Weight Formula''보다 ``통합된 Framework, Theory, Extensions, Cross-domain Validation''에 두는 것이 안전함.
\end{tcolorbox}

\section{Section 3 Proposed Method에 대한 평가}

\subsection{Framework라는 제목과 세 방법의 관계}

현재 논문 제목은 ``Framework''를 사용하며, Section 3은 LWCE, PLWCE, ES-LWCE를 하나의 Logarithmic Re-weighting Family로 설명함. 이러한 구성은 제목과 잘 연결됨.

다만 Framework라는 표현을 더욱 설득력 있게 하려면 세 방법을 개별적으로 나열하기 전에 일반적인 Weight Mapping을 먼저 정의하는 것도 고려할 수 있음.

예를 들어

\[
w_c
=
\psi\!\left(\log(1+n_c)\right),
\]

또는

\[
w_c(\alpha,\epsilon)
=
\frac{1}{\left(\log(1+n_c)+\epsilon\right)^\alpha}
\]

를 General Family로 제시한 뒤,

\begin{itemize}
\item LWCE: \((\alpha,\epsilon)=(1,0)\)
\item PLWCE: \(\alpha>0,\epsilon=0\)
\item ES-LWCE: \(\alpha=1,\epsilon>0\)
\end{itemize}

로 설명할 수 있음.

그러나 현재 원고는 Combined Variant를 Main Method로 제안하지 않으려는 방향이므로, General Formula를 전면에 제시하면 Reviewer가 ``왜 Combined를 사용하지 않는가''를 질문할 수 있음. 따라서 다음 두 방식 중 하나를 선택해야 함.

\begin{enumerate}
\item 현재처럼 세 방법을 분리하되 ``family''는 공통 Log Compression에서 나온다는 의미로만 사용함.
\item General Formula를 제시하고, \(\alpha\)와 \(\epsilon\)의 역할이 실제로 Orthogonal하지 않을 수 있음을 인정한 뒤 Ablation으로 선택을 정당화함.
\end{enumerate}

현재 Story에는 첫 번째 방식이 더 안전함.

\subsection{LWCE Weight의 절대 Scale}

LWCE의 Raw Weight

\[
w_c^{\mathrm{LWCE}}
=
\frac{1}{\log(1+n_c)}
\]

는 \(n_c>1\)일 때 대부분 \(1\)보다 작음. 따라서 Mean Normalization 없이 사용하면 Loss Scale을 감소시키는 문제가 있음. 현재 원고는 Mean Normalization을 적용하므로 이 문제를 해결하고 있음.

다만 이 Normalization 이후 Individual Weight의 Boundedness는 Raw Weight의 Bound와 동일하지 않음. Normalized Weight는 Class Composition에 따라 달라질 수 있음.

따라서 Proposition 2의 Boundedness는 Raw Weight에 대한 것임을 계속 명확히 유지하고, 실제 Training에서 사용되는 \(\widetilde{w}_c\)에 대해서는 별도의 Bound를 주장하지 않는 것이 좋음.

\subsection{PLWCE에서 \(\alpha\)의 해석}

현재 원고는

\[
w_c^{\mathrm{PLWCE}}
=
\frac{1}{\left(\log(1+n_c)\right)^\alpha}
\]

를 사용하고, \(\alpha<1\), \(\alpha=1\), \(\alpha>1\)로 해석함.

이 해석은 전반적으로 맞지만, Mean Normalization 이후의 실제 Tail Emphasis가 단순히 Raw Weight의 크기만으로 결정되는 것은 아님. 중요한 것은 Class 간 Relative Ratio임.

\[
\frac{w_{\min}^{\mathrm{PLWCE}}}{w_{\max}^{\mathrm{PLWCE}}}
=
\left(
\frac{\log(1+n_{\max})}{\log(1+n_{\min})}
\right)^\alpha.
\]

따라서 \(\alpha\)의 효과를 설명할 때 이 Ratio를 함께 제시하는 것이 좋음.

\subsection{ES-LWCE에서 \(\epsilon\)의 해석}

현재 원고는 \(\epsilon\)이 Weight를 Cap하고 Smoothing한다고 설명함. 그러나 Mean Normalization을 적용하면 Raw Weight의 Absolute Cap \(1/\epsilon\)이 실제 Normalized Weight의 동일한 Cap을 의미하지는 않음.

따라서 ``explicit cap on the raw per-class weight before normalization''이라고 표현하는 것이 더 정확함.

또한 \(\epsilon\to\infty\)일 때 Raw Weight가 모두 \(0\)으로 가지만, Mean Normalization 이후에는 Uniform Weight로 수렴함. 따라서 ``closer to uniform CE''라는 설명은 Mean-normalized Weight를 전제로 할 때 맞음. 이 전제를 명시하면 좋음.

\subsection{Proposition 1: Monotonicity}

Monotonicity는 간결하고 명확함. 다만 Class Count는 이산형이므로 Derivative 대신 Log Function의 Strict Monotonicity를 이용한 현재 Proof가 적절함.

추가로 Mean Normalization 후에도 Class Ordering이 보존됨을 Corollary로 한 문장 추가할 수 있음.

\begin{quote}
Because mean normalization multiplies all class weights by the same positive constant, it preserves the ordering of class weights.
\end{quote}

\subsection{Proposition 2: Boundedness}

현재 Proposition은 Raw Weight에 대해 정확함.

\[
0<w_c^{\mathrm{LWCE}}\leq \frac{1}{\log 2},
\qquad
0<w_c^{\mathrm{PLWCE}}\leq \frac{1}{(\log 2)^\alpha}
\]

for \(n_c\geq1\), 그리고

\[
0<w_c^{\mathrm{ES}}\leq \frac{1}{\epsilon}
\]

for \(n_c\geq0\)임.

다만 PLWCE의 경우 \(\alpha\)가 매우 클 때 \(1/(\log2)^\alpha\)가 크게 증가할 수 있음. 따라서 PLWCE는 \(n_c\)에 대해서는 bounded하지만, \(\alpha\)에 대해 Uniformly Bounded한 것은 아님.

추천 문장은 다음과 같음.

\begin{quote}
For any fixed \(\alpha>0\), PLWCE is bounded over \(n_c\geq1\).
\end{quote}

\subsection{Proposition 3: Dynamic-range Compression}

Proposition 3은 논문의 핵심 이론 결과임. WCE의 Linear Ratio와 LWCE의 Logarithmic Ratio를 직접 비교한 점이 좋음.

다만 PLWCE의 \(\alpha>1\)에서는 \(O((\log\rho)^\alpha)\)이며, \(\alpha\)가 충분히 크면 실질적인 Dynamic Range가 다시 크게 증가할 수 있음. 따라서 ``the proposed family grows only logarithmically''라는 포괄적 문장은 PLWCE에 대해서는 정확하지 않음.

LWCE는 Logarithmic, PLWCE는 Polylogarithmic, ES-LWCE는 Logarithmic보다 더 완화된 형태라고 구분해야 함.

\begin{tcolorbox}[title=수정 권장]
Abstract, Highlights, Introduction, Proposition 3, Conclusion에서 ``the proposed family grows only logarithmically''라고 일반화하지 말고, ``logarithmically for LWCE and polylogarithmically for PLWCE''라고 구분할 것.
\end{tcolorbox}

\subsection{Figure 1}

Figure 1은 WCE, Inverse-square-root, PLWCE, LWCE, ES-LWCE의 Weight Ratio를 비교하여 Dynamic-range Compression을 직관적으로 보여줌. 이론적 차이를 독자가 빠르게 이해할 수 있다는 점에서 유용함.

다만 다음 사항을 검토할 필요가 있음.

\begin{itemize}
\item \(\alpha=2,3\)인 PLWCE가 LWCE보다 훨씬 큰 Ratio를 보이므로, ``Compression Family가 항상 작은 Dynamic Range를 유지한다''는 인상을 주지 않도록 설명할 것.
\item ES-LWCE의 \(\epsilon=1\)은 실제 실험에서 사용하는 값인지, 단순 Illustration인지 Caption에 명시할 것.
\item Majority Count \(10^4\) 선택 이유를 설명할 것.
\item X-axis가 Minority Count이고 Ratio의 최대가 \(n_c=1\)에서 나타나는 구조를 본문에서 한 번 더 설명할 것.
\item Color만으로 Method를 구분하지 않고 Line Style과 Marker를 함께 사용하는 것이 좋음.
\end{itemize}

\subsection{Proposition 4: Gradient-scale Control}

Weighted CE의 Logit Gradient는 Sample \(i\)에 대해

\[
\frac{\partial \mathcal{L}_i}{\partial z_{i,k}}
=
w_{y_i}\left(p_{i,k}-\mathbb{I}[k=y_i]\right)
\]

의 형태임. 이 식을 Proposition 또는 Proof에 직접 포함하면 Class Weight가 Gradient에 Multiplicative Factor로 들어간다는 점이 더 분명해짐.

현재 ``full gradient magnitude also depends on predicted probabilities and feature representation''이라고 설명한 점은 적절함.

추가로 다음 두 가지를 구분하면 좋음.

\begin{itemize}
\item Logit-level Gradient Scaling
\item Parameter-level Gradient Norm
\end{itemize}

Logit-level에서는 \(w_{y_i}\)가 직접 곱해지지만, Network Parameter Gradient는 Jacobian과 Feature에 따라 달라짐. Experiment에서는 둘 중 무엇을 측정하는지 명확히 해야 함.

\section{Section 3.3 Parameter Selection의 추가 보완}

\subsection{Proxy Search의 공정성}

PLWCE만 Proxy Subset Search를 사용하고 다른 Tunable Baseline은 Full Validation Search를 사용하면 Search Budget이 달라질 수 있음. 현재 Section 4.2는 Tunable Baseline에도 동일한 Automated Budget을 부여한다고 설명하므로 방향은 좋음.

다만 ``동일한 Budget''을 다음 기준 중 무엇으로 정의하는지 명시해야 함.

\begin{itemize}
\item 동일 Candidate 수
\item 동일 Training Epoch 수
\item 동일 Proxy Subset 크기
\item 동일 총 GPU Time
\end{itemize}

가장 단순하고 재현 가능한 방식은 모든 Tunable Method에 동일 Proxy Subset과 동일 Candidate 수를 적용하는 것임.

\subsection{Proxy Subset의 대표성}

Severe Imbalance에서는 Stratified Proxy Subset을 구성할 때 Minority Class Sample이 지나치게 적어질 수 있음. ``minimum number of samples per class''를 보장한다고 했으므로, 다음 사항을 명시할 필요가 있음.

\begin{itemize}
\item 복원추출 여부
\item Minority Class Sample이 Minimum보다 적을 때 처리 방식
\item Proxy가 원래 Training Distribution을 유지하는지
\item Validation Metric의 Variance를 줄이기 위해 반복하는지
\end{itemize}

\subsection{선택 기준}

Macro-F1을 선택 기준으로 사용한 것은 적절함. 그러나 Main Metric이 Macro-F1과 Worst-class Accuracy 두 개라면, \(\alpha\)를 Macro-F1만으로 선택하는 이유를 설명하는 것이 좋음.

Worst-class Accuracy는 매우 불안정할 수 있으므로 Macro-F1을 Validation Metric으로 사용하고 Worst-class Accuracy는 평가 전용으로 두는 방식이 합리적임. 이를 한 문장 설명하면 좋음.

\section{Section 3.4 ES-LWCE Practical Use Case 작성 방향}

ES-LWCE의 필요성을 실제로 보여주기 위해서는 다음과 같은 별도 실험 시나리오가 필요함.

\subsection{Mini-batch Count Scenario}

Dataset-level Count 대신 각 Mini-batch에서 추정한 \(n_{c,b}\)를 사용함. 일부 Batch에서 특정 Class가 나타나지 않으면 \(n_{c,b}=0\)이 됨.

다음 방법을 비교할 수 있음.

\begin{itemize}
\item Batch-level LWCE without safeguard
\item Batch-level LWCE with ad hoc clipping
\item ES-LWCE
\item Dataset-level LWCE
\end{itemize}

측정 항목은 다음과 같음.

\begin{itemize}
\item NaN/Inf 발생 횟수
\item Maximum Weight
\item Gradient Norm
\item Training Loss Stability
\item Macro-F1
\end{itemize}

\subsection{Streaming or Drifting Count Scenario}

Training Stream에서 Class Prior가 시간에 따라 변화하도록 구성함. Moving Count 또는 Cumulative Count를 사용하고, 초기 구간에서 아직 관측되지 않은 Class를 포함함.

이 경우 ES-LWCE가 단순한 Benchmark Variant가 아니라 실제 Dynamic Count 상황을 위한 안정화 방법임을 보여줄 수 있음.

\subsection{Extreme Rare Class Scenario}

일부 Class의 Training Sample을 \(0,1,2,5,10\) 수준으로 조정하여 Weight와 성능을 비교함. 다만 \(n_c=0\)이면 해당 Class에 대한 Training Sample 자체가 없으므로, 단순 Supervised Training에서는 그 Class의 Loss가 계산되지 않음.

따라서 \(n_c=0\)의 의미를 정확히 설정해야 함. Dataset-level에서 Class가 완전히 없는 상황이 아니라, Mini-batch나 Online Estimation Window 안에서 일시적으로 Count가 \(0\)인 상황이어야 ES-LWCE의 의미가 있음.

\begin{tcolorbox}[title=중요]
``Class count가 0이어도 Loss를 계산할 수 있다''고 오해되지 않도록 할 것. 해당 Sample이 없으면 그 Class의 Loss Term도 없음. ES-LWCE의 장점은 Dynamic Frequency Table을 유지하거나 Batch-adaptive Weight를 계산할 때 Weight Function 자체가 항상 정의된다는 데 있음.
\end{tcolorbox}

\section{Section 4 Experiment의 권장 전체 구조}

현재 Section 4는 기본 Skeleton만 작성되어 있으므로, 다음과 같은 구조를 권장함.

\begin{enumerate}
\item Experimental Questions
\item Datasets and Imbalance Construction
\item Models and Training Protocol
\item Baselines
\item Hyperparameter Selection
\item Evaluation Metrics
\item Main Results
\item Performance by Imbalance Severity
\item Optimization Stability and Mechanism
\item Ablation Study
\item Sensitivity Analysis
\item ES-LWCE Practical Scenario
\item Computational Cost
\item Statistical Analysis
\end{enumerate}

\subsection{4.1 Datasets and Imbalance Construction}

각 Domain에 대해 다음 정보를 표로 정리하는 것이 좋음.

\begin{longtable}{p{3.0cm}p{2.4cm}p{2.4cm}p{2.4cm}p{4.0cm}}
\toprule
Dataset & Domain & Classes & Samples & Imbalance setting \\
\midrule
\endfirsthead
\toprule
Dataset & Domain & Classes & Samples & Imbalance setting \\
\midrule
\endhead
CIFAR-10-LT & Image & 10 & To be filled & Exponential / Step, \(\rho\) values \\
CIFAR-100-LT & Image & 100 & To be filled & Exponential / Step, \(\rho\) values \\
Tabular datasets & Tabular & To be filled & To be filled & Natural or induced imbalance \\
Intrusion datasets & Network & To be filled & To be filled & Natural severe imbalance \\
Emotion dataset & Text & To be filled & To be filled & Long-tailed construction \\
\bottomrule
\end{longtable}

Synthetic Long-tail을 만드는 경우 다음 식을 명시할 수 있음.

\[
n_c
=
n_{\max}\rho^{-\frac{c-1}{C-1}},
\qquad c=1,\ldots,C.
\]

또는 Step Imbalance라면 Head와 Tail Group의 Sample 수를 명확히 정의해야 함.

\subsection{4.2 Model and Training Protocol}

각 Domain별로 Model이 다르므로 다음을 명시해야 함.

\begin{itemize}
\item Architecture와 Layer 수
\item Initialization
\item Optimizer
\item Initial Learning Rate
\item Learning-rate Schedule
\item Batch Size
\item Epoch
\item Weight Decay
\item Data Augmentation
\item Early Stopping 여부
\item Random Seed
\item Hardware와 Software Version
\end{itemize}

논문의 목적은 Loss 비교이므로 Architecture와 Training Protocol은 모든 Loss에 동일해야 함.

\subsection{4.3 Baseline}

현재 Baseline은 CE, WCE, Inverse-square-root, CB Loss, Focal Loss, Logit Adjustment를 포함함. 기본 비교군으로 적절함.

추가로 다음 중 하나를 고려할 수 있음.

\begin{itemize}
\item LDAM 또는 LDAM-DRW
\item Balanced Softmax
\item Seesaw Loss
\item Equalization Loss
\end{itemize}

다만 모든 최신 Long-tailed Method를 포함할 필요는 없음. 본 논문은 Architecture나 Decoupled Training이 아니라 Simple Re-weighting Family를 다루므로, 직접 비교 가능한 Loss-level Baseline을 중심으로 구성하는 것이 더 중요함.

\begin{tcolorbox}[title=Baseline 공정성]
각 Baseline의 Hyperparameter를 Authors' Recommended Value만 사용하는 경우와 Validation으로 Tune하는 경우가 혼재하지 않도록 할 것. Parameter-free Method와 Tunable Method를 비교할 때 Tunable Method에는 합리적인 Validation Budget을 제공해야 함.
\end{tcolorbox}

\subsection{4.4 Evaluation Metric}

Primary Metric으로 Macro-F1과 Worst-class Accuracy를 사용하는 방향은 적절함. 다음과 같이 역할을 구분할 수 있음.

\begin{itemize}
\item Macro-F1: Class별 Precision과 Recall을 균형 있게 반영하는 Primary Metric
\item Worst-class Accuracy: 가장 취약한 Class의 성능을 직접 평가
\item Balanced Accuracy: Mean Per-class Recall
\item Geometric Mean: 일부 Class의 매우 낮은 Recall에 민감
\item Minority Recall: Tail Group의 평균 Recall
\item Overall Accuracy: 참고 지표
\end{itemize}

Worst-class Accuracy는 Seed에 따라 변동성이 클 수 있으므로 Mean \(\pm\) Standard Deviation과 Confidence Interval을 함께 보고하는 것이 좋음.

\subsection{4.5 Main Result Table}

Main Table은 Domain별로 분리하는 것이 좋음. 한 표에 모든 Dataset과 Metric을 넣으면 지나치게 복잡해질 수 있음.

권장 형식은 다음과 같음.

\begin{center}
\begin{tabular}{lcccc}
\toprule
Method & Macro-F1 & Worst Acc. & Balanced Acc. & G-Mean \\
\midrule
CE & & & & \\
WCE & & & & \\
Inv.-sqrt & & & & \\
Focal & & & & \\
CB Loss & & & & \\
Logit Adj. & & & & \\
LWCE & & & & \\
PLWCE & & & & \\
ES-LWCE & & & & \\
\bottomrule
\end{tabular}
\end{center}

Best는 Bold, Second-best는 Underline을 사용할 수 있음. 다만 Statistical Significance가 없는 작은 차이를 지나치게 강조하지 않는 것이 좋음.

\subsection{4.6 Imbalance Severity에 따른 성능 변화}

논문의 핵심 Claim은 Imbalance가 심해질수록 Logarithmic Compression의 장점이 커진다는 것임. 따라서 단일 Imbalance Ratio의 Table만으로는 부족함.

다음 그래프가 중요함.

\begin{itemize}
\item X-axis: Imbalance Ratio \(\rho\)
\item Y-axis: Macro-F1
\item Methods: CE, WCE, CB, Focal, LWCE, PLWCE
\end{itemize}

추가로

\[
\Delta_{\mathrm{LWCE-CE}}(\rho)
=
\mathrm{MacroF1}_{\mathrm{LWCE}}(\rho)
-
\mathrm{MacroF1}_{\mathrm{CE}}(\rho)
\]

를 표시하면 ``Imbalance가 커질수록 Gain이 증가하는가''를 직접 확인할 수 있음.

\section{Optimization Stability와 Mechanism 분석}

\subsection{Epoch별 Training Loss}

CE, WCE, LWCE, PLWCE의 Epoch별 Training Loss를 동일 Scale Convention에서 비교해야 함. Mean-normalized Weight를 사용하므로 Loss Scale 차이는 크게 줄어들겠지만, 여전히 곡선의 Smoothness와 Oscillation을 비교할 수 있음.

보고할 수 있는 정량 지표는 다음과 같음.

\begin{itemize}
\item Loss Curve의 Standard Deviation
\item Epoch-to-epoch Loss Change
\item Maximum Loss Spike
\item Divergence 또는 NaN 발생 여부
\end{itemize}

\subsection{Gradient Norm}

Gradient Norm은 다음 중 무엇을 측정했는지 명확히 해야 함.

\begin{itemize}
\item 전체 Parameter Gradient의 \(\ell_2\)-norm
\item Final Classifier Layer Gradient Norm
\item Per-sample Gradient Norm
\item Class별 Average Gradient Norm
\end{itemize}

논문의 Mechanism을 가장 직접적으로 보여주려면 Class별 또는 Head/Tail Group별 Gradient Norm을 계산하는 것이 좋음.

\[
G_c^{(t)}
=
\frac{1}{|\mathcal{B}_c|}
\sum_{i\in\mathcal{B}_c}
\left\|
\nabla_{\theta}\mathcal{L}_i^{(t)}
\right\|_2.
\]

그리고

\[
R_G^{(t)}
=
\frac{\frac{1}{|\mathcal{T}|}\sum_{c\in\mathcal{T}}G_c^{(t)}}
{\frac{1}{|\mathcal{H}|}\sum_{c\in\mathcal{H}}G_c^{(t)}}
\]

를 Tail-to-head Gradient Ratio로 사용할 수 있음.

단, Per-sample Gradient 계산은 비용이 크므로 Final Layer 또는 Representative Batch에서 측정할 수 있음.

\subsection{Class Weight Distribution}

각 Method의 Normalized Weight를 Class Frequency 순서로 Plot하면 WCE와 LWCE의 Dynamic Range 차이를 직접 보여줄 수 있음.

다음 값을 함께 보고하는 것이 좋음.

\begin{itemize}
\item \(\max_c\widetilde{w}_c\)
\item \(\min_c\widetilde{w}_c\)
\item \(\max_c\widetilde{w}_c/\min_c\widetilde{w}_c\)
\item Weight의 Coefficient of Variation
\end{itemize}

\subsection{Training Stability의 정의}

``Optimization Stability''는 다소 넓은 개념이므로 논문에서 Operational Definition을 제시하는 것이 좋음.

예를 들어 다음 중 일부를 사용함.

\begin{itemize}
\item Gradient Norm의 최대값
\item Gradient Norm의 분산
\item Loss Spike의 빈도
\item Seed 간 최종 성능 분산
\item Convergence Epoch
\item Divergence Rate
\end{itemize}

\begin{tcolorbox}[title=핵심]
WCE보다 LWCE의 Macro-F1이 높다는 사실만으로 Optimization Stability가 향상되었다고 결론내리면 안 됨. Stability Claim은 Gradient, Loss Dynamics, Seed Variance 등 직접적인 지표로 별도 입증해야 함.
\end{tcolorbox}

\section{Ablation Study 작성 방향}

현재 Table 4의 Ingredient Ablation은 구조가 좋음.

\begin{center}
\begin{tabular}{lccc}
\toprule
Method & Log Compression & Power \(\alpha\) & Smoothing \(\epsilon\) \\
\midrule
WCE & X & X & X \\
LWCE & O & X & X \\
PLWCE & O & O & X \\
ES-LWCE & O & X & O \\
Combined & O & O & O \\
\bottomrule
\end{tabular}
\end{center}

다만 Ablation은 단순히 Method별 성능을 나열하는 것보다 다음 질문을 직접 답하도록 구성해야 함.

\begin{enumerate}
\item WCE에서 LWCE로 바꾸었을 때 Log Compression 자체가 얼마나 기여하는가?
\item \(\alpha\)가 추가되면 어떤 Imbalance 수준에서 이득이 발생하는가?
\item \(\epsilon\)은 Fixed-count Benchmark에서 실제로 Accuracy를 높이는가?
\item Combined가 PLWCE보다 나은가?
\item Mean Normalization이 결과에 어떤 영향을 주는가?
\end{enumerate}

마지막 항목을 위해 다음 비교를 Supplementary에 넣을 수 있음.

\begin{itemize}
\item Raw Weight
\item Sum-normalized Weight
\item Mean-normalized Weight
\item Max-normalized Weight
\end{itemize}

이 비교는 Main Method를 변경하기 위한 것이 아니라 Scale Convention의 Robustness를 확인하기 위한 것임.

\subsection{Combined Variant의 Negative Result}

현재 원고는 Combined Variant가 PLWCE보다 개선되지 않고, \(\epsilon\)이 Grid의 Lower End로 선택된다고 미리 서술하고 있음. 실제 결과가 아직 확정되지 않았다면 이러한 문장은 결과가 나온 뒤 작성해야 함.

또한 ``consistently drives \(\epsilon\) to the lower end''라는 표현은 Dataset별 Selected Value Table로 뒷받침해야 함.

Negative Result를 유지하려면 다음을 제시하는 것이 좋음.

\begin{itemize}
\item Dataset별 선택된 \((\alpha,\epsilon)\)
\item PLWCE와 Combined의 성능 차이
\item Search Cost 차이
\item \(\epsilon\)이 거의 \(0\)으로 선택되는 빈도
\end{itemize}

\section{Sensitivity Analysis 작성 방향}

\subsection{\(\alpha\) Sensitivity}

\(\alpha\) Candidate는 예를 들어

\[
\alpha\in\{0.25,0.5,0.75,1,1.5,2,3\}
\]

처럼 설정할 수 있음. 단, 실제 Grid는 충분한 Preliminary Experiment 후 확정해야 함.

다음 Plot을 권장함.

\begin{itemize}
\item \(\alpha\) vs. Macro-F1
\item \(\alpha\) vs. Worst-class Accuracy
\item \(\alpha\) vs. Tail Recall
\item \(\alpha\) vs. Maximum Gradient Norm
\end{itemize}

Imbalance Ratio별 Curve를 함께 제시하면 \(\rho\)가 클수록 Optimal \(\alpha\)가 증가하는지 확인할 수 있음.

\subsection{\(\epsilon\) Sensitivity}

\(\epsilon\)은 Log Scale Grid를 사용하는 것이 적절함.

\[
\epsilon\in\{10^{-4},10^{-3},10^{-2},10^{-1},1\}.
\]

다만 \(\epsilon\)의 의미는 Weight Formula의 Log Base와 Mean Normalization에 따라 달라지므로, 단순 숫자만 제시하지 말고 그 값이 유도하는 Raw Cap \(1/\epsilon\)도 함께 표기하면 좋음.

\subsection{Hyperparameter Robustness}

PLWCE가 Tunable Method이므로 최적값에서만 우수한 것이 아니라 넓은 \(\alpha\) 범위에서 안정적인지 보여주는 것이 중요함.

다음 지표를 사용할 수 있음.

\[
\mathrm{RobustnessRange}
=
\left\{
\alpha:
\mathrm{MacroF1}(\alpha)
\geq
\max_{\alpha'}\mathrm{MacroF1}(\alpha')-\tau
\right\}.
\]

여기서 \(\tau\)는 허용 성능 차이임.

\section{Statistical Analysis 작성 방향}

Dataset과 Imbalance Level이 여러 개이면 각 Dataset--Imbalance 조합을 하나의 Block으로 간주하여 Method를 비교할 수 있음.

권장 절차는 다음과 같음.

\begin{enumerate}
\item 각 Block에서 Method Rank 계산
\item Friedman Test
\item 유의한 경우 Holm-corrected Pairwise Post-hoc Test
\item Mean Rank 보고
\item Win/Tie/Loss Count 보고
\end{enumerate}

단, Seed별 결과를 독립 Dataset처럼 취급하면 안 됨. Seed는 같은 Experimental Condition의 반복이므로, 먼저 Seed 평균을 구한 뒤 Condition 단위로 Rank를 계산하는 것이 일반적으로 더 적절함.

\subsection{Metric별 검정}

Primary Metric이 두 개이면 Macro-F1과 Worst-class Accuracy 각각 별도 검정을 수행할 수 있음. 다만 Multiple Metric에 대한 추가 Multiplicity를 고려해야 함.

가장 단순한 방식은 Macro-F1을 Primary Statistical Endpoint로 지정하고, Worst-class Accuracy는 Supporting Analysis로 두는 것임.

\subsection{Effect Size}

\(p\)-value만 제시하지 말고 다음을 함께 보고하면 좋음.

\begin{itemize}
\item Mean Rank Difference
\item Median Performance Difference
\item Win/Tie/Loss
\item Pairwise Effect Size
\end{itemize}

\section{Abstract 작성 방향}

현재 Abstract는 전반적으로 논문의 문제, 방법, 확장, Domain, 결과의 흐름을 갖추고 있음. 다만 실제 Experiment 결과가 아직 채워지지 않았으므로 최종 단계에서 다음 원칙에 따라 다시 작성하는 것이 좋음.

\begin{enumerate}
\item Problem: WCE의 Relative Weight Dynamic Range가 Severe Imbalance에서 과도하게 증가함.
\item Gap: Simplicity를 유지하면서 이를 직접 완화하는 Frequency-based Method가 부족함.
\item Method: LWCE, PLWCE, ES-LWCE.
\item Theory: Bounded Raw Weight와 Log/Polylog Ratio Growth.
\item Experiment: Domain 수, Dataset 수, Imbalance 범위.
\item Result: 실제 수치 또는 일관된 Win Count.
\item Conclusion: Simple, Stable, Parameter-efficient Alternative.
\end{enumerate}

\begin{tcolorbox}[title=Abstract에서 수정할 표현]
``inverse-frequency weights grow without bound''는 Ratio 또는 Normalized Weight의 의미로 정밀화할 것. ``avoid collapse'' 또는 ``severe degradation''은 실제 결과가 충분히 일관될 때만 유지할 것. Domain 수는 본문과 일치시킬 것.
\end{tcolorbox}

\section{Highlights 수정 방향}

현재 Highlights는 논문의 핵심을 비교적 잘 요약함. 다만 다음 표현을 조정하는 것이 좋음.

\begin{itemize}
\item ``bounds the weight explosion of WCE''는 ``compresses the rare-to-frequent weight dynamic range''로 구체화함.
\item ``single-parameter generalization tuned automatically''에서 Proxy Search Cost를 과장 없이 표현함.
\item ``keeps per-class weights bounded even for the rarest classes''는 Raw Weight와 \(n_c=0\) 조건을 구분함.
\item ``avoid collapse of rival losses''는 경쟁 방법을 지나치게 부정적으로 표현할 수 있으므로 완화함.
\end{itemize}

예시는 다음과 같음.

\begin{itemize}
\item Logarithmic compression reduces the dynamic range of inverse-frequency class weights.
\item LWCE is parameter-free and preserves stronger emphasis on rare classes.
\item PLWCE tunes tail emphasis through a single exponent selected on training-only proxy data.
\item ES-LWCE remains well defined for vanishing or batch-adaptive class counts.
\item Cross-domain experiments evaluate performance and optimization stability under increasing imbalance.
\end{itemize}

\section{Conclusion 작성 방향}

Conclusion은 단순한 Abstract 반복이 아니라 다음 내용을 포함하는 것이 좋음.

\begin{enumerate}
\item 핵심 문제와 해결 원리
\item 주요 이론 결과
\item 주요 실험 결과
\item Practical Advantage
\item Limitation
\item Future Work
\end{enumerate}

\subsection{반드시 포함할 Limitation}

\begin{itemize}
\item LWCE가 모든 Dataset과 Imbalance 수준에서 항상 최선이라고 주장할 수 없음.
\item PLWCE는 \(\alpha\) Selection Cost가 있음.
\item ES-LWCE의 장점은 Dynamic Count Scenario에서 더 직접적이며, Fixed Dataset-level Count에서는 효과가 제한될 수 있음.
\item Class Count만 사용하는 Frequency-based Method이므로 Sample Difficulty나 Label Noise를 직접 모델링하지 않음.
\item Extreme Tail에서 Rare Class Sample 자체가 매우 적으면 Re-weighting만으로 Representation 부족을 해결하기 어려움.
\item Large-scale Foundation Model Fine-tuning에 대한 검증이 제한적일 수 있음.
\end{itemize}

\subsection{Future Work}

\begin{itemize}
\item Difficulty-based Modulation과 Log Weight의 결합
\item Dynamic 또는 Online Class Count Estimation
\item Multi-label Classification
\item Open-set Long-tailed Learning
\item Label Noise와 Class Imbalance의 동시 처리
\item Foundation Model Fine-tuning
\item Theoretically Optimal Compression Function 탐색
\end{itemize}

\section{논문 전체에서 용어를 통일할 것}

다음 용어를 일관되게 사용하는 것이 좋음.

\begin{longtable}{p{5cm}p{9cm}}
\toprule
권장 용어 & 사용 원칙 \\
\midrule
Class imbalance & 일반적인 불균형 문제 \\
Long-tailed classification & Head-to-tail 구조가 명확한 설정 \\
Class count \(n_c\) & Frequency와 혼용하지 않고 수식에서는 Count 사용 \\
Raw weight \(w_c\) & Normalization 이전 Weight \\
Normalized weight \(\widetilde{w}_c\) & 실제 Loss에 사용하는 Weight \\
Weight dynamic range & Rare-to-frequent Ratio를 의미할 때 사용 \\
Weight explosion & 반드시 정의 후 제한적으로 사용 \\
Logarithmic compression & LWCE의 핵심 Mechanism \\
Polylogarithmic growth & PLWCE Ratio를 설명할 때 사용 \\
Optimization stability & 측정 지표를 정의한 뒤 사용 \\
Tail emphasis & Minority Class에 대한 Relative Weight 강조 \\
\bottomrule
\end{longtable}

\section{논리적 일관성 점검}

최종 원고에서는 다음 Claim--Evidence 연결이 명확해야 함.

\begin{longtable}{p{6cm}p{8cm}}
\toprule
Claim & Required evidence \\
\midrule
LWCE controls WCE weight explosion & Weight ratio theory and weight distribution plot \\
LWCE improves optimization stability & Gradient norm, loss curve, seed variance \\
LWCE helps more under severe imbalance & Performance across multiple \(\rho\) values \\
PLWCE controls tail emphasis & \(\alpha\) sensitivity and weight-ratio analysis \\
ES-LWCE handles vanishing counts & Mini-batch or streaming experiment \\
The method is domain-independent & Consistent results across all claimed domains \\
The method is parameter-efficient & Hyperparameter count and search-cost comparison \\
The method is computationally light & Runtime and memory comparison \\
\bottomrule
\end{longtable}

\begin{tcolorbox}[title=핵심 점검]
Introduction에서 제시한 Claim이 Experiment에서 직접 검증되지 않으면 해당 Claim을 완화해야 함. 반대로 Experiment에서 중요한 결과가 확인되면 Abstract, Contribution, Conclusion에 반영해야 함. 각 Section이 독립적으로 작성된 것처럼 보이지 않도록 Claim과 Evidence를 끝까지 연결하는 것이 중요함.
\end{tcolorbox}

\section{현재 원고의 주요 장점}

\begin{enumerate}
\item 문제 설정이 단순하고 직관적임.
\item WCE의 Dynamic-range 문제와 제안 방법의 관계가 명확함.
\item LWCE가 Parameter-free라는 Practical Advantage가 있음.
\item PLWCE와 ES-LWCE의 목적을 서로 다르게 정의한 점이 좋음.
\item Theory가 단순한 Formula 제안보다 한 단계 확장되어 있음.
\item Image, Tabular, Intrusion, Text로 확장 가능한 Domain-independent Loss임.
\item Architecture를 변경하지 않는 Drop-in Method라는 장점이 있음.
\item Weight Normalization을 통해 Baseline 비교의 공정성을 고려함.
\end{enumerate}

\section{현재 원고에서 가장 우려되는 점}

\subsection{Novelty}

Logarithmic Class Weight 또는 Inverse-log Weight가 기존에 사용된 사례가 있을 가능성이 있음. 따라서 Literature Search를 충분히 하지 않으면 Novelty Claim이 약해질 수 있음.

\subsection{Weight Explosion 용어}

Raw WCE Weight \(1/n_c\) 자체의 Unboundedness와 Rare-to-frequent Ratio의 증가를 혼동하면 Reviewer가 수학적 문제를 지적할 가능성이 있음.

\subsection{ES-LWCE의 실험적 필요성}

Dynamic Count Experiment가 없으면 ES-LWCE가 단순한 추가 Variant처럼 보일 수 있음.

\subsection{Strong Baseline}

Long-tailed Learning 분야의 강한 Loss Baseline이 부족하다고 평가될 수 있음. 최소한 Logit Adjustment 외에 Balanced Softmax 또는 LDAM 계열 중 하나를 추가하는 것을 검토하는 것이 좋음.

\subsection{Generality Claim}

네 Domain을 모두 언급하는 만큼 모든 Domain에서 충분한 Dataset과 공정한 Tuning을 제공해야 함. Domain별 Dataset이 하나뿐이면 Generality Claim을 약간 완화하는 것이 좋음.

\section{수정 우선순위}

\subsection{Priority 1: 반드시 먼저 완료할 사항}

\begin{enumerate}
\item Section 4의 Dataset, Model, Training Protocol을 정확한 수치로 작성할 것.
\item 본문 전체에서 평가 Domain 수를 통일할 것.
\item Weight Explosion을 Rare-to-frequent Dynamic Range 중심으로 정밀하게 정의할 것.
\item Main Result를 여러 Imbalance Ratio에서 제시할 것.
\item Gradient Norm, Loss Curve, Class Weight Distribution을 포함할 것.
\item ES-LWCE를 위한 Mini-batch 또는 Dynamic Count Experiment를 설계할 것.
\item PLWCE 및 모든 Tunable Baseline의 공정한 Validation Budget을 확정할 것.
\item Novelty와 관련된 Logarithmic Class Weight 선행연구를 추가 조사할 것.
\end{enumerate}

\subsection{Priority 2: 논문의 설득력을 크게 높이는 사항}

\begin{enumerate}
\item Balanced Softmax 또는 LDAM 계열 Baseline 추가를 검토할 것.
\item \(\alpha\), \(\epsilon\) Sensitivity Analysis를 작성할 것.
\item Combined Variant와 Mean Normalization Ablation을 수행할 것.
\item Runtime, Memory, Search Cost를 비교할 것.
\item Friedman Test, Holm Procedure, Mean Rank, Win/Tie/Loss를 보고할 것.
\item Domain별 Per-class 또는 Head/Medium/Tail 성능을 제시할 것.
\item Selected \(\alpha\)와 \(\epsilon\) 값을 Dataset별로 표로 정리할 것.
\end{enumerate}

\subsection{Priority 3: 최종 완성도 향상}

\begin{enumerate}
\item Abstract와 Highlights를 실제 결과에 맞게 다시 작성할 것.
\item Related Work의 Taxonomy 용어를 일관되게 정리할 것.
\item Figure 1의 Caption과 시각적 구분을 개선할 것.
\item Limitations와 Future Work를 구체적으로 작성할 것.
\item Supplementary Material에 Hyperparameter Grid, Detailed Results, Additional Curves를 포함할 것.
\item 코드 공개가 가능하면 Reproducibility Statement를 추가할 것.
\end{enumerate}

\section{권장 최종 Section 구성}

\begin{enumerate}
\item Introduction
\item Related Work
\begin{enumerate}
\item Taxonomy of Imbalanced Learning
\item Frequency-based Re-weighting
\item Difficulty-based and Margin-based Objectives
\item Representation and Architecture-level Methods
\item Positioning of Logarithmic Re-weighting
\end{enumerate}
\item Proposed Method
\begin{enumerate}
\item Logarithmic Re-weighting Family
\item Theoretical Properties
\item PLWCE Parameter Selection
\item ES-LWCE for Dynamic Counts
\item Complexity and Implementation
\end{enumerate}
\item Experiments
\begin{enumerate}
\item Research Questions
\item Datasets and Imbalance Construction
\item Models and Training Protocol
\item Baselines and Tuning
\item Evaluation Metrics
\item Main Results
\item Effect of Imbalance Severity
\item Optimization Stability
\item Ablation and Sensitivity
\item Dynamic-count Evaluation
\item Computational Cost
\item Statistical Analysis
\end{enumerate}
\item Discussion
\begin{enumerate}
\item When Logarithmic Compression Helps
\item Trade-off between Tail Recall and Head Accuracy
\item Practical Recommendations
\item Limitations
\end{enumerate}
\item Conclusion
\end{enumerate}

\section{최종 종합 의견}

현재 원고는 3차 피드백에서 요구한 수학적 정밀화와 Method Positioning을 상당히 잘 반영하였음. 특히 Raw Weight Boundedness와 Weight Ratio Compression을 구분하고, Gradient-scale Claim을 제한하며, Mean Normalization과 Parameter Selection 절차를 명시한 점은 중요한 개선임.

지금부터 가장 중요한 것은 새로운 Formula나 Method를 추가하는 것이 아니라, 이미 제시한 Claim을 실험으로 정확히 검증하는 것임. 특히 다음의 세 축이 끝까지 유지되어야 함.

\begin{center}
Controlled Weight Dynamic Range

$\Downarrow$

Controlled Class-dependent Gradient Amplification

$\Downarrow$

Reliable Performance under Severe Imbalance
\end{center}

첫 번째 축은 Proposition과 Figure로 비교적 잘 설명되어 있음. 두 번째 축은 Gradient와 Optimization 분석이 필요하며, 세 번째 축은 여러 Domain과 Imbalance Ratio에서의 Main Result 및 Statistical Analysis가 필요함.

\begin{tcolorbox}[title=최종 판단]
Introduction부터 Section 3까지는 3차 피드백을 전반적으로 잘 반영했으며, 논문의 기본 구조와 핵심 메시지는 안정된 상태임. 이후 작성에서는 성능표만 채우는 방식보다, Weight Compression이 실제 Gradient Dynamics와 Tail Performance에 어떤 영향을 주는지를 단계적으로 보여주는 것이 중요함. 이 Mechanism이 명확하게 검증되면 단순한 Re-weighting Variant가 아니라, Class Imbalance에서 Weight Dynamic Range를 통제하는 일반적인 Logarithmic Re-weighting Framework로 설득력 있게 Positioning할 수 있음.
\end{tcolorbox}

\end{document}